% Generated draft: requires mathematical and source review.
\documentclass[11pt]{article}
\usepackage[margin=25mm]{geometry}
\usepackage{fontspec}
\setmainfont{Latin Modern Roman}
\usepackage{amsmath,amssymb,mathtools,unicode-math}
\setmathfont{Latin Modern Math}
\usepackage{enumitem,needspace,xurl,hyperref}
\setlength{\parindent}{0pt}
\setlength{\parskip}{5pt}
\setlength{\emergencystretch}{3em}
\newcommand{\N}{\mathbb{N}}
\newcommand{\Z}{\mathbb{Z}}
\newcommand{\Q}{\mathbb{Q}}
\newcommand{\R}{\mathbb{R}}
\newcommand{\C}{\mathbb{C}}
\newcommand{\F}{\mathbb{F}}
\newcommand{\A}{\mathbb{A}}
\newcommand{\PP}{\mathbb{P}}
\newcommand{\E}{\mathbb{E}}
\newcommand{\eps}{\varepsilon}
\newcommand{\dd}{\,\mathrm{d}}
\begin{document}
\section*{UnsolvedMath ID: 10800004}
\textbf{Source problem code:} catalogue entry
\subsection{Definitions and mathematical statement}
\textbf{Problem (Victor Vassiliev, Problem 1D).} Give more general lower bounds for the dimension of $\mu=\mathrm{const}$ strata in terms of the intersection forms of vanishing cycles.

In the customary isolated complex hypersurface singularity interpretation, a holomorphic germ $f:(\C^n,0)\to(\C,0)$ with an isolated critical point has Milnor number
\[
 \mu(f)=\dim_{\C}\frac{\mathcal O_{\C^n,0}}{(\partial f/\partial z_1,\ldots,\partial f/\partial z_n)}.
\]
Its Milnor fiber $F_f$ has a vanishing-cycle lattice $H_{n-1}(F_f;\Z)$ equipped with its integral intersection pairing $I_f$. In a specified deformation space with base $B$, let $\Sigma_\mu$ denote the germ of the locus where the relevant deformed singularity has Milnor number $\mu(f)$. The requested quantitative bounds have the form
\[
 \dim_{\C}\Sigma_\mu\ \ge\ b(I_f),
\]
where the lower bound is obtained from the intersection-form data. This display describes the type of bound sought; the source does not propose a particular function $b$ or assert a particular numerical inequality.

The supplied excerpt does not fix the deformation space, the convention for following the critical point, the permitted singularity classes or dimensions, or whether the input consists of the pairing up to isomorphism or a distinguished vanishing-cycle configuration. These choices must be recovered from the surrounding source. In particular, arbitrary additional trivial deformation parameters change stratum dimension, so a convention such as a specified miniversal deformation matters. The registry summary suggests counting ``independent prohibited collisions,'' but supplies neither the collision restrictions nor a definition of independence. It also reports a positivity/non-simplicity relationship; that report alone is not a precise general quantitative bound.
\subsection{Short English statement}
Find quantitative lower bounds for the dimension of a constant-Milnor-number stratum using the intersection pairing of its vanishing cycles. The singularity and deformation conventions, and any independence notion for prohibited collisions, still need to be specified.
\subsection{Sources}
\begin{itemize}
\item Local UnsolvedMath registry, permanent ID 10800004, source code catalogue entry; proposer Victor Vassiliev, no proposal year supplied.
\item Victor Vassiliev, \emph{A Few Problems on Monodromy and Discriminants}, Problem 1D; \url{https://amj.math.stonybrook.edu/html-articles/Files-2015-2024/15-11/}. The supplied provenance records HTML-structure extraction, access on 2026-07-29, and eighteen explicitly labeled theorem-style problems in the article.
\item The registry research summary reports positivity of the constant-$\mu$ stratum dimension in relation to non-simplicity and discusses independent prohibited collisions. It gives no separate primary citation or precise hypotheses for those reports.
\end{itemize}
\end{document}
